\hypertarget{appendix-2a2e5bd7143a}{}\label{referencestart:appendix-2a2e5bd7143a}
\section{Mastery evidence levels}\label{appendix-2a2e5bd7143a-mastery-evidence-levels}

\hypertarget{appendix-2a2e5bd7143a-mastery-evidence-levels}{}

PROJECTALLMIGHT should not use one word such as ``developed'' for fundamentally different evidence states. A concept can be well explained but not executable, executable in a CPU reference but not on the native GPU, or validated in a teaching lab but not in the current production repository.

Use the following evidence ladder.

\subsection{0 --- Indexed}\label{appendix-2a2e5bd7143a-indexed}

\hypertarget{appendix-2a2e5bd7143a-indexed}{}

The topic has a real node, canonical name, prerequisites and destination. This proves discoverability only.

\subsection{1 --- Explained}\label{appendix-2a2e5bd7143a-explained}

\hypertarget{appendix-2a2e5bd7143a-explained}{}

The forcing problem, definitions, conventions and causal mental model are stated without relying on unexplained jargon.

\subsection{2 --- Derived and traced}\label{appendix-2a2e5bd7143a-derived-and-traced}

\hypertarget{appendix-2a2e5bd7143a-derived-and-traced}{}

The central mathematics or algorithm is derived with symbol meanings, assumptions and at least one nontrivial worked trace containing intermediate state.

\subsection{3 --- Reference implemented}\label{appendix-2a2e5bd7143a-reference-implemented}

\hypertarget{appendix-2a2e5bd7143a-reference-implemented}{}

A complete, rebuildable implementation exists in the teaching surface or a linked standalone laboratory. Positive and negative cases execute.

\subsection{4 --- Independently verified}\label{appendix-2a2e5bd7143a-independently-verified}

\hypertarget{appendix-2a2e5bd7143a-independently-verified}{}

The mechanism is checked against an oracle, invariant, property test, differential implementation, numerical bound or external normative reference strong enough to expose plausible bugs.

\subsection{5 --- Native-system verified}\label{appendix-2a2e5bd7143a-native-system-verified}

\hypertarget{appendix-2a2e5bd7143a-native-system-verified}{}

Claims involving a real API/device/tool are exercised on that system: Vulkan validation/device output, Metal GPU capture, Blender/Maya/Houdini/Nuke host execution, real filesystem semantics, real network stack, and so on.

\subsection{6 --- Production-repository grounded}\label{appendix-2a2e5bd7143a-production-repository-grounded}

\hypertarget{appendix-2a2e5bd7143a-production-repository-grounded}{}

The current Kairo/Maveb/Vulkax/Nanoquant or other target repository revision is inspected and the exact implementation, ownership, lifetime, synchronization, performance and failure paths are tied back to the lesson.

\subsection{7 --- Reconstructed mastery}\label{appendix-2a2e5bd7143a-reconstructed-mastery}

\hypertarget{appendix-2a2e5bd7143a-reconstructed-mastery}{}

The learner can derive, implement, trace, profile, debug and compare alternatives without copying the source, and can solve a transfer problem in a different context.

A concept can hold several evidence dimensions separately. For example, Vulkan synchronization may be Derived=Yes, CPU reference=Yes, Native Vulkan=No, Kairo parity=Unverified. Do not collapse that into one optimistic ``complete'' flag.

\subsection{Required proof fields}\label{appendix-2a2e5bd7143a-required-proof-fields}

\hypertarget{appendix-2a2e5bd7143a-required-proof-fields}{}

For each M001--M180 chunk, the evidence ledger should record:

\begin{itemize}
\tightlist
\item
  canonical section;
\item
  original part range;
\item
  definitions and conventions;
\item
  derivation;
\item
  worked trace;
\item
  executable/reference artifact;
\item
  positive/negative/property tests;
\item
  alternatives/failure modes;
\item
  external source review;
\item
  native platform evidence if applicable;
\item
  production repository evidence;
\item
  research bridge;
\item
  retrieval/oral reconstruction status.
\end{itemize}

The purpose is not bureaucracy. It prevents a three-megabyte note from masking one missing proof obligation.
